
\section{Discussion}\label{sec:discussion}

\paragraph{What the evidence supports.} On the real data we examined, simple methods are hard
to beat. On pairs the best method on the chambers is RECI, which compares two regression errors, and on
Tübingen no method is distinguishable from RECI, with LOCI nominally ahead; on graphs the
best methods on recordings with independent rows are DirectLiNGAM and GES, and on the Sachs
data a mix of methods that use nonparametric tests, with differences smaller than the sampling
error of a $17$-edge reference. The methods we developed to address specific failures show their
intended benefit in simulation and some of it on measurements: lower false-edge rates under
repeated inspection and nonlinear mechanisms in simulation, and the highest precision on the
chamber recordings, BTAD, at the price of recall. None of them has the best $F_1$ on the
chambers. The two classifiers trained on simulated data are the clearest negative results, and
we think they are informative: a model that is nearly perfect on a broad simulator can be at
chance on real pairs.

\paragraph{Implications for evaluation.} First, the effective sample size of a recording
should be reported: seventeen of nineteen new chamber recordings had fewer than twenty
effective observations, which no test that assumes independent rows can cope with, and which
also contaminates reference sets built with such tests. Second, adjusting the score for
dependences the reference graph cannot explain changes conclusions and should be shown next to the
strict score, with the construction of the adjustment disclosed and made independent of the methods.
Third, the protocol matters: several of our early positive results were produced by
choices made after seeing the data, and the frozen evaluation reversed them. We report which
numbers are frozen and which are not (Section~\ref{sec:provenance}).

\paragraph{Limitations.} The reference networks are small, the Sachs consensus is contested,
and the chamber graphs are bipartite, so recovery of the skeleton is mostly feature selection
from actuators to measurements. Most differences between graph methods are within sampling
error, and the effective sample size of the walk recordings limits what any method can show.
We did not evaluate amortised methods (AVICI), CAM or QCCD, and several baselines
are re-implementations or default runs (LOCI on at most $1000$ rows); a better-tuned baseline could change a ranking.
The Causal Chambers also contain interventional experiments, where a simple shift test already
recovers the descendants with perfect precision, and we did not use them. The new methods were
developed by the author of this paper, which is a source of bias; the frozen protocol and the
public code are the mitigation.
